% !TeX root = ../main.tex

\chapter*{前言}
\addcontentsline{toc}{chapter}{前言}

这本专著对具有分数维数与整数维数 Hausdorff 维数的集合的数学理论给出了严谨而自足的论述.它主要关注几何理论而非应用.书中的许多内容此前只能见于原始数学论文，而这些论文往往高度专门、令人困惑，并且使用陈旧的记号.我希望通过撰写本书使这些材料更易于获取，同时也为使用分形集合的读者提供一份有用而精确的论述.

本书虽然主要为纯数学家而写，但我希望它对若干类型或多或少有素养且有较高要求的读者也有用处.在最基本的层次上，本书可供在其他数学或科学学科中遇到分形的读者作为参考书使用：将主要定理与推论同基本定义结合起来阅读，即可得到那些已被严格建立的性质的精确表述.

要想获得这一学科的全貌，或者在初次阅读时，可以只跟随基本的解说与定理的陈述而略去详细的证明.非专业的数学工作者或许还可以略去~1.1~节的细节，该节从技术的观点建立了一般测度的性质.

要充分领会书中的细节，需要具备初等数学分析与一般拓扑的实用知识.毫无疑问，构成这一理论核心的某些证明既困难又冗长，但为了完整洞察其中数学的美妙与精巧，掌握它们是十分值得的.

本书着重于该学科的基本工具，例如 Vitali 覆盖定理、网测度与位势论方法.

我们所需要的测度与 Hausdorff 测度的性质在第 1 章的前两节中建立.全书着重于把测度就其自身用于估计集合的大小，而不是将其作为定义积分的一个步骤.积分只在后面各章中作为一个方便的工具使用；大体上，关于积分的直观认识就完全够用了.

在所采用的普遍性水平上不可避免地做了一些折中.我们在 $n$ 维 Euclid 空间中工作，尽管许多思想同样适用于更一般的度量空间.在某些情形，高维类似物的证明要复杂得多，此时定理只在二维情形给出证明，并对其推广给出参考文献.类似地，有时也只给出一维或二维的证明，如果它们包含一般情形的本质思想，同时又能使记号得到简化.我们还把注意力限制在与数值维数 $s$ 相应的 Hausdorff 测度上，而不是任意函数对应的 Hausdorff 测度.

书中不少证明相对原始形式做了简化.进一步的缩短无疑也是可能的，但作者对完美性的追求不得不与在有限时间内完成本书的要求相权衡！

虽然本专著本质上是自足的，书中还是简要描述了这一工作的各种变形与推广，并给出了完整的参考文献.更多的变形与推广可以在各章末尾的习题中找到.

在此谨向所有以任何方式帮助过本专著写作的人们致谢.我特别感谢 Roy Davies 教授对书稿的细致批评，并允许我使用其未发表的资料；感谢 Hallard Croft 博士提出的详尽建议以及在校对中给予的帮助.我也非常感谢 B.~B.~Mandelbrot 教授、J.~M.~Marstrand 教授、P.~Mattila 教授和 C.~A.~Rogers 教授有益的意见与讨论.

我还要感谢 Maureen Woodward 夫人与 Rhoda Rees 夫人为书稿打字，并感谢 Cambridge University Press 的 David Tranah 与 Sheila Shepherd 陪伴本书走过了出版的各个阶段.最后，我必须感谢我的妻子 Isobel 拨冗阅读了本书的早期草稿，以及她始终如一的鼓励与支持.
